\documentclass[10pt,a4paper]{amsart}
\usepackage{amsmath,amssymb,amsthm,graphicx}
% Retained minimal preamble for the Stacks Project chapter more-algebra.tex (More on Algebra).
% The chapter file pinned-original-source.tex carries only \input{preamble} on its first line, so no
% preamble text lives inside the chapter itself.  The Stacks Project uses one single shared preamble for
% every chapter; the snapshot stages it identically as sources/stacks-chow/preamble.tex and
% sources/stacks-etale-cohomology/preamble.tex.  The block below is copied mechanically, line for line,
% from sources/stacks-chow/preamble.tex lines 14-18 (the Xy-pic setup) and 154-217 (the native theorem
% declarations together with the mathematical macros).  Nothing was added, removed or redefined; the
% wrapper class and the amsmath/amssymb/amsthm/graphicx base packages are supplied by the assembler.
% Chapter: more-algebra.tex -- More on Algebra.
% For commutative diagrams we use Xy-pic
\usepackage[all]{xy}

% We use 2cell for 2-commutative diagrams.
\xyoption{2cell}

% Theorem environments.
%
\theoremstyle{plain}
\newtheorem{theorem}[subsection]{Theorem}
\newtheorem{proposition}[subsection]{Proposition}
\newtheorem{lemma}[subsection]{Lemma}

\theoremstyle{definition}
\newtheorem{definition}[subsection]{Definition}
\newtheorem{example}[subsection]{Example}
\newtheorem{exercise}[subsection]{Exercise}
\newtheorem{situation}[subsection]{Situation}

\theoremstyle{remark}
\newtheorem{remark}[subsection]{Remark}
\newtheorem{remarks}[subsection]{Remarks}

\numberwithin{equation}{subsection}

% Macros
%
\def\lim{\mathop{\mathrm{lim}}\nolimits}
\def\colim{\mathop{\mathrm{colim}}\nolimits}
\def\Spec{\mathop{\mathrm{Spec}}}
\def\Hom{\mathop{\mathrm{Hom}}\nolimits}
\def\Ext{\mathop{\mathrm{Ext}}\nolimits}
\def\SheafHom{\mathop{\mathcal{H}\!\mathit{om}}\nolimits}
\def\SheafExt{\mathop{\mathcal{E}\!\mathit{xt}}\nolimits}
\def\Sch{\mathit{Sch}}
\def\Mor{\mathop{\mathrm{Mor}}\nolimits}
\def\Ob{\mathop{\mathrm{Ob}}\nolimits}
\def\Sh{\mathop{\mathit{Sh}}\nolimits}
\def\NL{\mathop{N\!L}\nolimits}
\def\CH{\mathop{\mathrm{CH}}\nolimits}
\def\proetale{{pro\text{-}\acute{e}tale}}
\def\etale{{\acute{e}tale}}
\def\QCoh{\mathit{QCoh}}
\def\Ker{\mathop{\mathrm{Ker}}}
\def\Im{\mathop{\mathrm{Im}}}
\def\Coker{\mathop{\mathrm{Coker}}}
\def\Coim{\mathop{\mathrm{Coim}}}

% Boxtimes
%
\DeclareMathSymbol{\boxtimes}{\mathbin}{AMSa}{"02}

%
% Macros for moduli stacks/spaces
%
\def\QCohstack{\mathcal{QC}\!\mathit{oh}}
\def\Cohstack{\mathcal{C}\!\mathit{oh}}
\def\Spacesstack{\mathcal{S}\!\mathit{paces}}
\def\Quotfunctor{\mathrm{Quot}}
\def\Hilbfunctor{\mathrm{Hilb}}
\def\Curvesstack{\mathcal{C}\!\mathit{urves}}
\def\Polarizedstack{\mathcal{P}\!\mathit{olarized}}
\def\Complexesstack{\mathcal{C}\!\mathit{omplexes}}
% \Pic is the operator that assigns to X its picard group, usage \Pic(X)
% \Picardstack_{X/B} denotes the Picard stack of X over B
% \Picardfunctor_{X/B} denotes the Picard functor of X over B
\def\Pic{\mathop{\mathrm{Pic}}\nolimits}
\def\Picardstack{\mathcal{P}\!\mathit{ic}}
\def\Picardfunctor{\mathrm{Pic}}
\def\Deformationcategory{\mathcal{D}\!\mathit{ef}}

% --- appended 2026-10-12 (Chief): the project own definitions of the reference/slogan/history
% display environments, copied mechanically from the shared Stacks preamble lines 8-12
% (\usepackage{verbatim} plus the three \newenvironment lines, which the chapter preamble keeps
% verbatim from the project preamble). Nothing added, removed or redefined: these environments
% carry no mathematical content and the project itself typesets them as comments. Without these
% three lines a counted statement whose own text uses them fails with Environment slogan undefined.
\usepackage{verbatim}
\newenvironment{reference}{\comment}{\endcomment}
\newenvironment{slogan}{\comment}{\endcomment}
\newenvironment{history}{\comment}{\endcomment}
\usepackage{cleveref}
\crefname{originallemma}{original lemma}{original lemmas}
\Crefname{originallemma}{Original lemma}{Original lemmas}
\makeatletter
\newlabel{algebra-lemma-characterize-unramified}{{02FM}{source}{}{sourcepointer-1}{}}
\newlabel{algebra-lemma-characterize-unramified@cref}{{[originallemma][0][]02FM}{[1][0][]source}{}{}{}}
\makeatother
\crefname{originallemma}{original lemma}{original lemmas}
\Crefname{originallemma}{Original lemma}{Original lemmas}
\makeatletter
\newlabel{algebra-lemma-completion-flat}{{00MB}{source}{}{sourcepointer-2}{}}
\newlabel{algebra-lemma-completion-flat@cref}{{[originallemma][0][]00MB}{[1][0][]source}{}{}{}}
\makeatother
\crefname{originallemma}{original lemma}{original lemmas}
\Crefname{originallemma}{Original lemma}{Original lemmas}
\makeatletter
\newlabel{algebra-proposition-unramified-locally-standard}{{0395}{source}{}{sourcepointer-3}{}}
\newlabel{algebra-proposition-unramified-locally-standard@cref}{{[originallemma][0][]0395}{[1][0][]source}{}{}{}}
\makeatother
\crefname{originallemma}{original lemma}{original lemmas}
\Crefname{originallemma}{Original lemma}{Original lemmas}
\makeatletter
\newlabel{lemma-henselization-noetherian}{{10979}{source}{}{sourcepointer-4}{}}
\newlabel{lemma-henselization-noetherian@cref}{{[originallemma][0][]10979}{[1][0][]source}{}{}{}}
\makeatother
\crefname{originallemma}{original lemma}{original lemmas}
\Crefname{originallemma}{Original lemma}{Original lemmas}
\makeatletter
\newlabel{lemma-quotient-by-idempotent}{{32037}{source}{}{sourcepointer-5}{}}
\newlabel{lemma-quotient-by-idempotent@cref}{{[originallemma][0][]32037}{[1][0][]source}{}{}{}}
\makeatother
\title{Every one-dimensional minimal prime of the completion of a Noetherian local ring is the radical of the extension of a minimal prime of the henselization}
\author{The Stacks Project Authors}
\date{stacks-project:more-algebra.tex:local-snapshot-20261012; GNU FDL 1.2 or later}
\begin{document}
\maketitle
\setcounter{subsection}{1}
\noindent\textit{Editorial note.} All delivered mathematical text is the author's own, copied from the pinned Stacks Project chapter file more-algebra.tex; no statement or proof step was rewritten, reordered or removed.\par\medskip
\noindent\textit{Editorial note.} The chapter file pinned-original-source.tex carries only \textbackslash{}input\{preamble\} on its first line, so no preamble text lives inside the chapter. The retained file sources/stacks-more-algebra/retained-preamble.tex is a mechanical, line-for-line copy of the project's single shared preamble as staged in the snapshot; nothing was added, removed or redefined and no package was added.\par\medskip
\noindent\textit{Editorial note.} The source URL is the source-level Stacks Project URL, as required when more than one card is produced from one source; the per-claim tag is carried by claim\_key.\par\medskip
\noindent\textit{Editorial note.} The proof is carried out with the auxiliary prime p = q cap A introduced only near the end; that order is the author's and is reproduced. The five references inside the counted proof span are recorded as non-load-bearing navigation pointers.\par\medskip
\noindent\textit{Allowed standard prerequisites (editorial).}\begin{itemize}
\item Henselization and completion of a Noetherian local ring, and their agreement. used at proof lines 32109-32110. Is the reduction that lets A be assumed henselian.
\item The quotient-by-idempotent lemma for a quotient of an adic completion killing a square-zero annihilated ideal. invoked at proof lines 32111-32120. Produces the finite type A-algebra C whose completion is A-hat/J.
\item The structure of unramified ring maps as quotients of etale algebras, and the rigidity of etale algebras over a henselian local ring with trivial residue extension. used at proof lines 32126-32133. Forces B = A after localization and hence C = A/I.
\item Exactness and flatness of the completion map. used at proof lines 32134-32137. Gives J = IA-hat and I = J cap A, which turn the containment into the radical statement.
\end{itemize}\medskip
\ifdefined\hypertarget\hypertarget{sourcepointer-1}{}\fi
\noindent\textit{Source reference.} Original lemma 02FM: focused-seeds/AGP/sources/stacks-snapshot-20261011/tags.txt, source lines 3197--3197. Cross-chapter navigation pointer into the Stacks Project; the standard fact it supplies is applied inside the counted proof span and the argument resting on it stays there.\par\medskip
\ifdefined\hypertarget\hypertarget{sourcepointer-2}{}\fi
\noindent\textit{Source reference.} Original lemma 00MB: focused-seeds/AGP/sources/stacks-snapshot-20261011/tags.txt, source lines 817--817. Cross-chapter navigation pointer into the Stacks Project; the standard fact it supplies is applied inside the counted proof span and the argument resting on it stays there.\par\medskip
\ifdefined\hypertarget\hypertarget{sourcepointer-3}{}\fi
\noindent\textit{Source reference.} Original lemma 0395: focused-seeds/AGP/sources/stacks-snapshot-20261011/tags.txt, source lines 4260--4260. Cross-chapter navigation pointer into the Stacks Project; the standard fact it supplies is applied inside the counted proof span and the argument resting on it stays there.\par\medskip
\ifdefined\hypertarget\hypertarget{sourcepointer-4}{}\fi
\noindent\textit{Source reference.} Original lemma 10979: sources/stacks-more-algebra/pinned-original-source.tex, source lines 10979--10979. Same-chapter navigation pointer; the statement it names is standard for the section and the argument that consumes it stays inside the counted proof span.\par\medskip
\ifdefined\hypertarget\hypertarget{sourcepointer-5}{}\fi
\noindent\textit{Source reference.} Original lemma 32037: sources/stacks-more-algebra/pinned-original-source.tex, source lines 32037--32037. Same-chapter navigation pointer; the statement it names is standard for the section and the argument that consumes it stays inside the counted proof span.\par\medskip
\begin{lemma}
\label{lemma-one-dimensional-formal-branch}
Let $(A, \mathfrak m)$ be a Noetherian local ring with henselization $A^h$.
Let $\mathfrak q \subset A^\wedge$ be a minimal prime with
$\dim(A^\wedge/\mathfrak q) = 1$. Then there exists a minimal
prime $\mathfrak q^h$ of $A^h$ such that
$\mathfrak q = \sqrt{\mathfrak q^hA^\wedge}$.
\end{lemma}

\begin{proof}
Since the completion of $A$ and $A^h$ are the same, we may assume
that $A$ is henselian (Lemma \ref{lemma-henselization-noetherian}).
We will apply Lemma \ref{lemma-quotient-by-idempotent}
to $A^\wedge \to A^\wedge/J$ where
$J = \Ker(A^\wedge \to (A^\wedge)_{\mathfrak q})$.
Since $\dim((A^\wedge)_\mathfrak q) = 0$ we see that
$\mathfrak q^n \subset J$ for some $n$. Hence $J/J^2$ is
annihilated by $\mathfrak q^n$. On the other hand $(J/J^2)_\mathfrak q = 0$
because $J_\mathfrak q = 0$. Hence $\mathfrak m$ is the only
associated prime of $J/J^2$ and we find that a power
of $\mathfrak m$ annihilates $J/J^2$. Thus the lemma applies
and we find that $A^\wedge/J = C^\wedge$ for some finite type
$A$-algebra $C$.

\medskip\noindent
Then $C/\mathfrak m C = A/\mathfrak m$ because $A^\wedge/J$ has the same
property. Hence $\mathfrak m_C = \mathfrak m C$ is a maximal ideal
and $A \to C$ is unramified at $\mathfrak m_C$
(Algebra, Lemma \ref{algebra-lemma-characterize-unramified}).
After replacing $C$ by a principal localization we may
assume that $C$ is a quotient of an \'etale $A$-algebra $B$, see
Algebra, Proposition \ref{algebra-proposition-unramified-locally-standard}.
However, since the residue field extension of $A \to C_{\mathfrak m_C}$
is trivial and $A$ is henselian, we conclude that $B = A$
again after a localization.
Thus $C = A/I$ for some ideal $I \subset A$ and it follows that
$J = IA^\wedge$ (because completion is exact in our situation by
Algebra, Lemma \ref{algebra-lemma-completion-flat}) and $I = J \cap A$
(by flatness of $A \to A^\wedge$). Since
$\mathfrak q^n \subset J \subset \mathfrak q$ we see that
$\mathfrak p = \mathfrak q \cap A$ satisfies
$\mathfrak p^n \subset I \subset \mathfrak p$.
Then $\sqrt{\mathfrak p A^\wedge} = \mathfrak q$ and the proof is complete.
\end{proof}

\end{document}
